\chapter{无线 IoT 系统}
\label{chap:wireless-iot}

无线 IoT 系统同时受射频链路、功耗、网络拓扑、安全身份、云端状态和设备生命周期约束。选择无线协议不能只比较标称速率和距离，还要评估频谱环境、节点密度、业务时延、待机时间、网关依赖、认证要求和量产测试能力。

一条“设备能够联网”的演示链路，与能够长期交付的产品仍有明显距离。后者必须回答：弱网下是否会形成重试风暴，长期离线后如何收敛状态，重复命令是否安全，凭据怎样轮换，用户转让设备时哪些身份应保留，以及量产线能否稳定识别射频缺陷。

\section{从无线工作负载选择协议}

协议选择应从业务工作负载开始，而不是从芯片功能列表开始。至少需要量化：

\begin{itemize}
  \item 单次上下行载荷、周期、突发量、最大可接受时延和丢失容忍度；
  \item 节点数量、移动性、覆盖范围、遮挡、部署密度和干扰环境；
  \item 设备是持续在线、周期唤醒，还是长期离线后批量同步；
  \item 是否依赖手机、接入点、边界路由器、运营商或专用网关；
  \item 电池容量、峰值电流能力、目标寿命和允许的维护周期；
  \item 所需的组网、发现、设备管理、互操作和法规认证生态。
\end{itemize}

\begin{table}[htbp]
  \centering
  \caption{常见无线技术的工程侧重点}
  \small
  \begin{tabular}{p{18mm}p{28mm}p{25mm}p{24mm}p{25mm}}
    \hline
    技术 & 典型拓扑 & 主要优势 & 主要约束 & 常见用途 \\
    \hline
    Wi-Fi & 基础设施网络 & 吞吐高、IP 原生 & 峰值功耗和配网 & 摄像头、网关、家电 \\
    BLE & 星形、广播、mesh & 低功耗、手机生态 & 载荷与连接参数 & 传感器、穿戴、配置 \\
    Thread & IPv6 mesh & 低功耗、自组织 & 需边界路由器 & 楼宇和智能家居 \\
    Zigbee & 应用层 mesh & 生态成熟 & 非 IP 应用模型 & 传统家居与照明 \\
    LoRaWAN & 星中之星 & 远距离、低功耗 & 低速率和占空比 & 表计、农业、遥测 \\
    蜂窝物联网 & 运营商网络 & 广域和移动性 & 模组、资费、认证 & 资产追踪、远程设备 \\
    \hline
  \end{tabular}
\end{table}

Matter 是基于 IP 的应用层互操作体系，可运行在 Wi-Fi、Thread 和以太网上，并通常使用 BLE 完成初始配网\cite{matter-spec}。它不能替代底层无线链路设计，也不会自动解决天线、共存、互联网接入和设备生命周期问题。

\section{链路预算、天线与共存}

\subsection{从接收功率到链路裕量}

链路预算由发射功率、天线增益、路径损耗和其他损耗共同决定：
\begin{equation}
  P_{rx}=P_{tx}+G_{tx}+G_{rx}-L_{path}-L_{misc}
\end{equation}

若接收机在选定 PHY、带宽和误包率条件下的灵敏度为 $S_{rx}$，链路裕量可近似写为：
\begin{equation}
  M_{link}=P_{rx}-S_{rx}
\end{equation}

正裕量只表示静态估算可接收，不代表产品可靠。人体遮挡、多径衰落、天线失谐、制造偏差、温度、电源噪声和同频干扰都会消耗裕量。设计评审应明确采用的最坏姿态、目标误包率和保留裕量，而不能仅以“实验室内 RSSI 大于某值”作为覆盖结论。

RSSI 通常是芯片内部估计量，其偏差和滤波方式依实现而异。不同芯片、PHY 或固件版本的 RSSI 不宜直接比较。覆盖验收应同时观察丢包率、重传、有效吞吐和业务完成时间。

\subsection{天线必须在整机中验证}

PCB 天线必须遵守参考设计的板厚、介质、馈线阻抗、净空和地平面要求。外壳、电池、屏幕、排线、金属紧固件和人体都会改变谐振与方向图，最终产品必须进行传导与辐射测试，不能只验证裸板。

样机阶段至少应保留匹配网络、可测量的射频路径和明确的天线版本。若量产中更换外壳材料、电池供应商或线缆走向，应触发射频影响评估，而不是把它视为纯结构变更。

\subsection{共存是系统调度问题}

2.4~GHz Wi-Fi、BLE、Thread 和 Zigbee 共用频段。共存设计需要信道规划、收发仲裁、天线隔离和流量调度；双无线芯片若缺少硬件共存接口，软件高负载时可能出现明显丢包和延迟长尾。

共存验证不能只让两个无线接口“同时打开”，而要制造最坏重叠，例如 Wi-Fi 持续上行时保持 BLE 低间隔连接、Thread mesh 路由变化时执行大包下载，或 USB~3、DDR 和显示接口同时产生宽带噪声。应比较单无线基线与并发场景下的误包率、P99 时延、断连次数和能耗。

\section{连接恢复与协议边界}

\subsection{Wi-Fi 重连与漫游}

Wi-Fi 在线状态至少包含扫描、认证、关联、地址获取、DNS、传输连接、应用鉴权和业务同步多个阶段。只有获得 IP 地址不代表云服务可用；socket 连接成功也不代表应用会话已经恢复。

重连状态机应区分 AP 不可见、口令错误、DHCP 超时、DNS 失败、证书失败、服务限流和本地时间不可信。永久配置错误不应无限快速重试，暂时故障则应采用带 jitter 的指数退避。设备还应在 AP 重启、DHCP 租约变化、NAT 映射失效和系统 suspend/resume 后重新验证整条链路。

移动设备若需要漫游，应实测扫描中断、切换阈值和业务时延。802.11k/v/r 等能力取决于终端与网络共同支持，应用不能假定所有现场 AP 都具备快速漫游。粘滞连接、同名不同网络和门户认证也应进入异常用例。

\subsection{BLE 广播、连接与绑定}

BLE 广播间隔影响发现时延、碰撞概率和功耗；扫描窗口与扫描间隔则影响中心设备的发现能力。广播载荷应只暴露发现所需信息，不应把长期身份或敏感状态直接作为可跟踪的明文标识。

建立连接后，connection interval、peripheral latency 和 supervision timeout 必须联合配置。过短间隔增加射频和唤醒开销，过长间隔会扩大控制时延；supervision timeout 还必须满足协议规定的参数关系。GATT MTU、Data Length 和链路层 PHY 的协商结果均可能低于本地最大值，应用分片不能依赖单一手机或操作系统版本。

断连重连后，应用应重新确认服务发现、通知订阅和授权状态。bonding 记录也有容量上限和淘汰策略；恢复出厂、手机丢失和所有权转移时，双方遗留的绑定信息都可能阻止重新连接。

\subsection{Thread 与 Matter 的责任边界}

Thread 提供低功耗 IPv6 mesh 网络，边界路由器负责在 Thread 网络与相邻 IP 网络之间转发流量。边界路由器不是业务云代理，互联网故障与 mesh 内部连通性也不是同一状态。

验证应覆盖 leader 变化、路由器掉电、网络分区与合并、边界路由器切换、网络数据集更新以及 sleepy end device 长时间休眠。节点恢复后不能假定原父节点、地址或路由仍然有效。

Matter 的 fabric 身份、访问控制和应用数据模型，与 Wi-Fi 口令或 Thread 网络数据集属于不同层。重新加入链路网络不等同于重新加入 Matter fabric；删除某个 fabric 也不一定要求销毁设备的出厂证明。产品必须把链路凭据、管理域身份和用户所有权分别建模。

\section{功耗预算覆盖完整业务周期}

平均电流由各状态电流与驻留时间加权：
\begin{equation}
  I_{avg}=\frac{\sum I_i t_i}{\sum t_i}
\end{equation}

无线发送往往只占很短时间，实际寿命可能由连接保持、频繁扫描、晶振启动、传感器预热和电源静态电流决定。必须用电流分析仪测量完整业务周期，而不是用芯片数据手册中的深睡电流直接估算整机寿命。

BLE connection interval、peripheral latency、Wi-Fi DTIM、Thread sleepy end device 轮询周期和蜂窝 PSM/eDRX 都是在延迟、可靠性和功耗之间取舍。云端心跳、DNS 刷新、证书握手、日志上报和 OTA 检查周期同样属于设备电源设计。

弱网会显著改变能耗模型。一次事务的能量应包含扫描、关联、握手、重传、等待和尾部保持时间。验收指标宜使用“每次成功业务事务的能量”和“规定弱网下的日均能量”，避免低成功率场景因少发数据而看似更省电。

\section{配网、设备身份与所有权}

配网是建立设备身份、网络凭据和管理权限的安全过程。产品应使用每设备唯一身份和一次性或受保护的 onboarding 信息，避免全产品共享默认口令。

推荐流程包括：
\begin{enumerate}
  \item 用户通过物理动作使设备进入有时限的配网状态；
  \item 管理端验证设备证明、出厂身份或 possession 信息；
  \item 双方建立认证加密通道；
  \item 在通道内传递最小必要的网络凭据和访问控制信息；
  \item 设备验证目标网络，原子持久化配置并报告最终结果；
  \item 配网窗口关闭，失败尝试具有速率限制和审计记录。
\end{enumerate}

“配网成功”应有明确提交点。若手机已显示成功而设备在持久化前掉电，双方必须能查询并收敛最终状态。设备不应长时间同时保留测试网络、旧用户网络和新用户网络的有效凭据。

所有权转移至少包含旧账户撤权、旧 fabric 或管理关系删除、用户数据清除、网络与应用凭据轮换、新所有者重新证明控制权。设备不可复制的出厂根身份通常可以保留，但旧用户授权不能随之保留。

产品宜区分网络重置、应用数据重置、恢复出厂和安全报废。恢复出厂设置必须明确删除哪些网络凭据、绑定记录、用户数据和设备密钥；安全报废还要撤销云端身份并处理不可擦除存储。设备根身份是否允许删除，取决于重新入网、返修和供应链追溯模型。

\section{MQTT 5 的业务语义}

MQTT 是面向发布/订阅的消息协议。MQTT~5 增加了会话过期、消息过期、原因码、流量限制和更多属性，但协议选项必须映射为业务契约，不能只采用库的默认值\cite{mqtt-v5}。

\begin{table}[htbp]
  \centering
  \caption{MQTT 机制与需要明确的业务问题}
  \small
  \begin{tabular}{p{30mm}p{100mm}}
    \hline
    机制 & 工程语义 \\
    \hline
    QoS 0 & 协议层尽力而为，断线或进程崩溃可丢失，适合可被新值覆盖的遥测 \\
    QoS 1 & 至少一次交付，接收方必须容忍 DUP 和业务重复执行 \\
    QoS 2 & MQTT 会话内协议交付一次，不自动保证数据库或物理动作 \mbox{exactly-once} \\
    Session Expiry & 决定断线后订阅和未完成 QoS 流程保留多久，需要与设备离线周期匹配 \\
    Retained Message & broker 保存某 topic 的最后保留值，适合快照，不是完整事件历史 \\
    Will Message & 异常断开后的延迟通知，只表示会话状态，不等价于设备物理故障 \\
    Message Expiry & 过期后不再转发，业务仍需校验命令自身的截止期和 generation \\
    Receive Maximum & 限制并行 QoS 1/2 流量，帮助把 broker 背压传递到受限设备 \\
    \hline
  \end{tabular}
\end{table}

Clean Start 与 Session Expiry 应联合设计。设备若每次都创建全新会话，离线期间消息和订阅状态可能丢失；若会话无限保留，broker 又会积累陈旧状态。客户端标识必须稳定且不在多台设备间复用，否则新连接可能顶掉旧连接。

topic 不是授权边界。broker 应依据设备身份限制可发布和可订阅的 topic，避免仅靠设备端代码自律。topic 命名还应区分命令、事件、状态快照和诊断数据，防止 retained 命令在设备每次重连时被再次执行。

QoS 不能替代业务确认。控制命令至少应携带稳定的 command ID、schema 版本、生成时间、过期时间和期望状态版本。设备应返回“已执行”“已执行过”“暂时失败”“永久拒绝”或“状态冲突”等可机器判断结果。

\begin{lstlisting}[language=C,caption={可重试 IoT 命令的抽象结构}]
struct cloud_command {
    uint64_t command_id;
    uint64_t target_generation;
    uint32_t schema_version;
    uint32_t expires_at;
    uint16_t operation, payload_length;
    uint8_t  payload[MAX_PAYLOAD];
};
\end{lstlisting}

设备应拒绝过期、版本不支持、generation 冲突和已执行的命令。不可安全重复的物理动作需要独立的 prepare/commit 或查询恢复机制，不能因使用 QoS~2 就宣称实现端到端 exactly-once。

\section{CoAP 与 LwM2M}

\subsection{CoAP 消息与资源模型}

CoAP 面向受限节点和受限网络，以请求/响应方式访问资源，常运行在 UDP 上\cite{rfc7252}。Confirmable 消息通过 ACK 和重传提供协议可靠性，Non-confirmable 消息不要求逐条确认；二者都不能替代业务层状态检查。

Token 用于关联请求与响应，Message ID 用于检测一段时间窗口内的重复消息，两者用途不同。重传 Confirmable 请求时，服务端必须避免让同一个写操作产生重复副作用。经代理或重启后，应用仍需要自己的幂等键和资源版本。

Observe 扩展允许客户端订阅资源变化，但观察关系是需要恢复的软状态\cite{rfc7641}。客户端重连或服务端重启后应重新注册，并根据序列或资源版本判断通知新旧。低频关键状态还应周期性校验，不能无限依赖一次订阅。

Block-Wise Transfer 使用 Block1 和 Block2 分块传输较大请求或响应\cite{rfc9177}。实现必须限制总对象大小、并发传输数、块缓存和超时，并在固件下载等场景校验完整对象摘要。分块成功只证明传输完成，不代表镜像已经被验证或安装。

\subsection{LwM2M 设备管理}

LwM2M 在 CoAP 之上定义设备管理和遥测语义，使用 Object、Object Instance、Resource 和 Resource Instance 组织数据\cite{oma-lwm2m}。采用标准对象有助于互操作，自定义对象则必须稳定分配 ID、权限、单位和版本。

典型生命周期包括：

\begin{enumerate}
  \item Bootstrap Server 下发或更新管理服务器、凭据和访问控制；
  \item Client 向 LwM2M Server 注册其 endpoint 与对象能力；
  \item Client 周期更新注册，服务器执行 Read、Write、Execute 或 Observe；
  \item 固件更新对象驱动下载、校验、安装、重启和结果上报状态机；
  \item 所有权转移或报废时撤销服务器关系和凭据。
\end{enumerate}

bootstrap 权限通常高于普通设备管理权限，地址和信任根不应由未认证网络输入覆盖。固件更新的下载完成、包验证、安装完成和试运行确认是不同状态，具体回滚与防回滚策略应与\mbox{第~\ref{chap:system-build-update}~章}一致。

\section{设备影子与云边状态收敛}

设备影子常把云端期望状态表示为 \texttt{desired}，设备实际确认状态表示为 \texttt{reported}。二者是用于弱网同步的状态模型，不应把“云端已写 desired”显示成“设备已经执行”。

每个字段都要定义权威方和合并规则。例如工作模式可以由云端设定，故障码只能由设备报告，本地按键修改目标温度时则需要决定本地优先、云端优先还是按 generation 比较。一个统一的 last-write-wins 时间戳通常不够，因为设备时钟可能不可信，网络也会重排消息。

可采用以下收敛规则：

\begin{itemize}
  \item 云端每次修改 desired 时增加单调 generation；
  \item 设备只应用比已确认 generation 更新且尚未过期的目标；
  \item reported 同时携带已应用 generation、结果码和设备观察值；
  \item 云端只有收到匹配 generation 的 reported 才显示“已生效”；
  \item 不允许远程修改的字段由设备权威，云端不得用旧快照覆盖；
  \item 冲突、拒绝和部分应用必须成为显式状态，而不是静默重试。
\end{itemize}

状态快照与事件流应分开。影子适合表达“当前应是什么”，审计事件用于回答“发生过什么”。只保留最后状态会丢失中间告警，只保存事件又会使长期离线设备恢复成本过高。

\section{弱网、长期离线与背压}

无线设备必须把断线视为正常状态。连接状态机应是可取消、可超时并能跨重启恢复的，不能由多个任务各自循环重连。指数退避应加入随机抖动和上限，服务端明确返回限流或重试等待时应优先遵守服务端信号。

离线队列必须有确定边界和数据策略：

\begin{description}
  \item[状态快照] 只保留每个 key 的最新值，旧值可合并；
  \item[周期遥测] 可降采样、聚合或按优先级丢弃，并记录丢失计数；
  \item[告警事件] 保留顺序、唯一 ID 和持久化确认，预留独立容量；
  \item[控制命令] 校验过期、generation 和幂等键，不因重连无条件重放；
  \item[审计记录] 需要防篡改和明确保留期，队列满时不能静默覆盖。
\end{description}

队列策略应覆盖掉电一致性、flash 磨损和 schema 升级。设备重新上线时不能瞬间倾倒全部积压数据；应接受服务端背压，限制在途请求，并让实时控制流量优先于历史遥测。通用的幂等、重试和背压契约见\mbox{第~\ref{chap:interface-protocol-data-contracts}~章}。

过期判断要选择正确时钟域。短期命令可使用接收后的单调时限；跨设备绝对截止期需要可信实时时间和允许误差。若时间尚未可信，应进入受限策略，而不是默认所有历史命令仍有效。

\section{TLS、DTLS 与对象安全}

基于 TCP 的 MQTT 或 HTTPS 通常使用 TLS~1.3\cite{rfc8446}；基于 UDP 的 CoAP 可使用 DTLS~1.3\cite{rfc9147}。握手需要证书或预共享密钥、随机数、算法协商和额外往返，内存、峰值 CPU、分片和会话恢复都必须在目标硬件上验证。

网关或代理终止 TLS/DTLS 后可以读取和修改应用明文。若 CoAP 消息需要跨中间节点保持端到端机密性、完整性和重放保护，可使用 OSCORE 对应用消息进行对象安全保护\cite{rfc8613}。OSCORE 不隐藏所有网络元数据，也不能替代端点授权和密钥生命周期管理。

\subsection{首次可信时间}

设备首次启动时常没有可信日历时间，但证书又包含有效期。直接关闭证书时间校验会把启动难题变成长期漏洞。可选方案包括：

\begin{itemize}
  \item 保存带回滚保护的可信时间下界，启动后只允许时间向前推进；
  \item 使用固化公钥或受限信任锚连接专用 bootstrap 服务，再获取认证时间；
  \item 在安全元件或受保护 RTC 中保持时间，并检测电池丢失或回拨；
  \item 时间未建立前只开放最小注册能力，不接受普通控制和更新操作。
\end{itemize}

NTP 响应本身若未认证，不能单独成为证书校验的信任来源。设备应记录时间质量和来源，绝对时间与单调超时的边界见\mbox{第~\ref{chap:embedded-network-time}~章}。

\subsection{凭据轮换与失效恢复}

证书生命周期应覆盖生产注入、首次注册、续期、吊销、所有权转移和设备报废。轮换协议要允许新旧信任材料在受控窗口内重叠，并采用“先验证新凭据、再原子切换、最后删除旧凭据”的顺序。

若续期只允许通过即将过期的证书完成，长期离线设备可能永久失联。产品应定义受限恢复凭据、现场维修或重新证明设备身份的路径。根信任锚轮换还需要双签名、交叉信任或等价过渡机制，不能把新根作为普通未认证配置下发。

TLS 只保护传输通道，不能替代固件签名、本地授权和安全启动。完整安全生命周期见\mbox{第~\ref{chap:security-lifecycle}~章}。

\section{可观测性与故障诊断}

无线日志应记录阶段和原因，而不是反复输出“连接失败”。建议为每次连接周期分配 correlation ID，并记录：

\begin{itemize}
  \item 扫描、认证、关联、地址获取、DNS、握手和应用登录各阶段耗时；
  \item 断连的本地状态、协议原因码、对端和网络类型；
  \item RSSI 或链路质量、重传率、PHY、信道和漫游次数；
  \item 队列深度、最老消息 age、丢弃/合并原因和在途数量；
  \item MQTT 会话恢复、重复消息、命令过期和 generation 冲突计数；
  \item 每次成功事务能量、唤醒原因以及异常长尾样本。
\end{itemize}

日志不得包含 Wi-Fi 口令、私钥、完整 token 或可用于长期跟踪用户的标识。高频链路事件需要采样和聚合，避免诊断流量反过来破坏无线性能。设备群指标和诊断包设计见\mbox{第~\ref{chap:observability-fleet}~章}。

\section{法规、认证与量产测试}

无线产品需要考虑目标市场的频谱法规、发射功率、带外辐射、接收机要求、SAR/暴露限制以及联盟认证。使用预认证模组可以降低风险，但天线、布局、外壳、时钟或射频参数变化仍可能要求评估或重新测试。

法规测试、联盟互操作测试和产线筛查目的不同。法规测试证明规定条件下满足强制要求；互操作认证验证协议行为；产线测试则以可重复、短节拍的方法识别装配和器件异常，不能彼此替代。

量产测试应覆盖晶振偏差、发射功率、接收路径、天线连接、设备身份和地址唯一性。单点 RSSI 容易受工装距离、屏蔽箱、线缆和黄金样机漂移影响，必须通过校准、限值样机和 GR\&R 验证测试系统。测试模式必须受控，交付固件不应留下未经授权即可进入的连续发射接口。

生产写入应把法规域、射频校准、MAC/EUI、设备证书和产品序列号绑定到同一追溯记录。烧录失败或返工时不得重复分配身份，也不能让测试凭据进入交付状态。

\section{验证矩阵与验收门}

无线验证需要覆盖空间、时间和系统负载，而不是只做一次端到端冒烟测试。

\begin{table}[htbp]
  \centering
  \caption{无线 IoT 验证矩阵示例}
  \small
  \begin{tabular}{p{27mm}p{48mm}p{55mm}}
    \hline
    维度 & 代表场景 & 主要观测量 \\
    \hline
    射频边界 & 不同距离、姿态、外壳、温度与同频干扰 & PER、重传、吞吐、P99 时延 \\
    网络故障 & AP/路由器重启、DHCP/DNS 失败、NAT 超时 & 恢复成功率、阶段耗时、重试负载 \\
    云端故障 & broker 切换、限流、证书轮换、服务降级 & 会话收敛、积压、重复和数据丢失 \\
    电源事件 & 连接、配网、队列提交和 OTA 中随机掉电 & 持久状态、身份、可恢复性 \\
    长期离线 & 队列接近上限、命令过期、schema 已升级 & 合并/淘汰、背压、升级兼容性 \\
    生命周期 & 转让、恢复出厂、返修、报废 & 旧权限撤销、数据清除、重新入网 \\
    \hline
  \end{tabular}
\end{table}

验收门应给出可重复的数字指标，例如在规定衰减和干扰下的业务成功率、断网恢复 P95/P99、24 小时离线后的收敛时间、最大积压不丢失的高优先级事件数，以及典型与弱网场景的每事务能量。只写“能够自动重连”无法形成回归门。

测试环境需要记录 AP、手机、边界路由器、broker、固件和射频配置版本。可编程衰减器、屏蔽箱、网络故障代理和电源控制器应纳入自动化，但最终仍需在真实部署场景验证传播、漫游和人群干扰。

\section{验证清单}

\begin{itemize}
  \item 协议选择是否基于业务周期、覆盖、节点密度、时延和电池预算？
  \item 天线是否在最终外壳、电池、线缆、姿态和温度状态下测量？
  \item 共存测试是否覆盖其他无线和高速接口的最坏并发负载？
  \item 断网、弱信号、AP 重启、网关迁移、服务限流和凭据失效能否恢复？
  \item MQTT、CoAP 或 LwM2M 的会话、重复、分块和过期语义是否映射到业务状态？
  \item 设备影子的 desired 与 reported 是否有权威方、generation 和冲突规则？
  \item 离线队列是否有容量、持久化、淘汰、背压和重新上线限速策略？
  \item 配网、所有权转移、恢复出厂、返修和报废是否具有明确安全边界？
  \item 首次可信时间、证书轮换和长期离线后的恢复路径是否经过演练？
  \item 功耗测量是否覆盖真实连接、握手、重传、积压同步和环境边界？
  \item 量产测试是否经过校准和 GR\&R，并能追溯射频参数与设备身份？
  \item 验收是否使用成功率、分位数、队列 age 和每事务能量等数字门槛？
\end{itemize}
